%! Author  = Omar Iskandarani
%! Date    = June 2026
%! ORCID   = 0009-0006-1686-3961
%! Target  = Physical Review A

% ============================================================
%  A Single-Postulate Closure for Localized Vortex-Core Models:
%  Deriving the Characteristic Circulation Speed from the
%  Rydberg Constant and the Compton Frequency
% ============================================================

\documentclass[aps,pra,preprint,longbibliography,nofootinbib]{revtex4-2}

\usepackage{amsmath,amssymb,amsthm,mathtools}
\usepackage{booktabs}
\usepackage{enumitem}
\usepackage{siunitx}
\usepackage[colorlinks=true,
            linkcolor=blue!60!black,
            citecolor=blue!60!black,
            urlcolor=blue!60!black]{hyperref}
\usepackage{cleveref}
\usepackage{microtype}

% ── Theorem environments ─────────────────────────────────────
\newtheorem{postulate}{Postulate}
\newtheorem{lemma}{Lemma}
\newtheorem{proposition}{Proposition}
\newtheorem{corollary}{Corollary}
\newtheorem{remark}{Remark}

% ── Notation ─────────────────────────────────────────────────
\newcommand{\vc}{v}              % characteristic circulation speed
\newcommand{\rc}{r_c}            % core radius
\newcommand{\omC}{\omega_C}      % Compton angular frequency
\newcommand{\lc}{\lambda_c}      % Compton wavelength
\newcommand{\Rinf}{R_\infty}     % Rydberg constant
\newcommand{\Fmax}{F_{\max}}     % maximum vortex tension
\newcommand{\rhocore}{\rho_{\mathrm{core}}} % core density
\newcommand{\Gz}{\Gamma_0}       % circulation quantum
% ─────────────────────────────────────────────────────────────

\begin{document}

\title{A Single-Postulate Closure for Localized Vortex-Core Models:\\
       Deriving the Characteristic Circulation Speed\\
       from the Rydberg Constant and the Compton Frequency}

\author{Omar Iskandarani}
\affiliation{Independent Researcher, Groningen, The Netherlands}
\email{info@omariskandarani.com}
\orcid{0009-0006-1686-3961}
\date{June 2026}

% ─────────────────────────────────────────────────────────────
\begin{abstract}
We present a minimal, self-consistent derivation of a characteristic
set of dynamical scales for localized vortex-core models of the electron
from a single structural postulate combined with the spectroscopic value
of the Rydberg constant.
The postulate, called the \emph{Compton-core hypothesis}, identifies
the angular rotation rate of the vortex core with the electron Compton
frequency $\omC = m_e c^2/\hbar$.
Together with the measured Rydberg constant $\Rinf$, this yields the
characteristic circulation speed
$\vc = \bigl(\pi\hbar c\Rinf/m_e\bigr)^{1/2}$
without invoking the fine-structure constant $\alpha$ as a prior input.
Instead, $\alpha = 2\vc/c$ emerges as a derived quantity and
reproduces the CODATA value to within $3\times 10^{-10}$.
The core radius $\rc = \vc/\omC$, the maximum vortex tension
$\Fmax = \hbar\omC/(2\rc)$, the core mass density
$\rhocore = m_e c^2/(2\pi\vc^2\rc^3)$, and the
circulation quantum $\Gz = 2\pi\vc^2/\omC$ are then derived
as algebraic consequences.
A status table distinguishes derived quantities from the single
calibrated input ($\Rinf$) and from open parameters.
Three falsifiable predictions are stated that do not involve
any free parameters.
\end{abstract}

\maketitle

% =============================================================
\section{Introduction}
\label{sec:intro}
% =============================================================

Localized vortex-filament models of fundamental particles trace their
conceptual lineage to the knotted-vortex proposals of Helmholtz and
Kelvin~\cite{Helmholtz1867,Kelvin1867}, and have received renewed
technical attention in the context of topological solitons~\cite{FaddeevNiemi1997},
analogue gravity~\cite{Barcelo2005}, and superfluid models of the
vacuum~\cite{Volovik2003}.
In each of these frameworks a characteristic circulation speed $\vc$
and a core radius $\rc$ appear as fundamental scales, but their
numerical values are typically supplied by hand or fitted to known
particle masses.

The present paper addresses a more basic question: can $\vc$ and $\rc$
be obtained from a minimal set of independently motivated constraints,
without circularity?

We show that the answer is yes, provided one accepts a single structural
postulate about the core's rotation rate.
Given that postulate and the spectroscopic value of the Rydberg
constant~\cite{MohrTaylor2018}, the two scales are uniquely determined,
and a hierarchy of further mechanical scales follows algebraically.
Notably, the fine-structure constant $\alpha$ is not used as an input;
it emerges as a consequence, and its derived value agrees with the
CODATA 2018 result~\cite{MohrTaylor2018} to ten significant figures.

The derivation is self-contained, dimensionally explicit at every step,
and numerically verified against CODATA values.
A status table (Sec.~\ref{sec:status}) labels each quantity as
\emph{derived}, \emph{calibrated}, or \emph{open} to make the logical
structure transparent to the reader.

Throughout we work in SI units.
$\hbar = h/(2\pi)$, $m_e$, $c$, and $e$ carry their standard CODATA
meanings.
The Rydberg constant $\Rinf$ is the only experimental input beyond
$\{m_e, c, \hbar\}$.

% =============================================================
\section{The Compton-Core Postulate}
\label{sec:postulate}
% =============================================================

The single structural assumption of this paper is:

\begin{postulate}[Compton-core hypothesis]
\label{post:compton}
The vortex core rotates at the electron Compton angular frequency,
\begin{equation}
  \Omega_{\mathrm{core}} = \omC \equiv \frac{m_e c^2}{\hbar}.
  \label{eq:postulate}
\end{equation}
The characteristic tangential speed at the core boundary of radius
$\rc$ is therefore
\begin{equation}
  \vc = \omC\,\rc.
  \label{eq:vc_rc_link}
\end{equation}
\end{postulate}

\begin{remark}
\label{rem:motivation}
In superfluid vortex physics the Compton frequency sets the natural
energy scale for the core: it is the frequency at which a quantum of
energy $\hbar\omC = m_e c^2$ is stored in the vortex
core~\cite{Donnelly1991,Pethick2008}.
Identifying the core rotation rate with $\omC$ is therefore the
simplest resonance condition consistent with the electron rest energy.
The postulate is testable in principle: any independent measurement
of $\rc$ inconsistent with $\rc = \vc/\omC$ would falsify it.
\end{remark}

\begin{remark}
Postulate~\ref{post:compton} introduces one equation relating two
unknowns, $\vc$ and $\rc$.
A second equation is needed to fix both scales uniquely.
We obtain it from the Rydberg constant in Sec.~\ref{sec:rydberg}.
\end{remark}

% =============================================================
\section{The Rydberg Constraint}
\label{sec:rydberg}
% =============================================================

\subsection{Statement}

The Rydberg constant is defined by atomic spectroscopy as
\begin{equation}
  \Rinf = \frac{\alpha^2 m_e c}{2h},
  \label{eq:Rydberg_standard}
\end{equation}
where $\alpha = e^2/(4\pi\varepsilon_0\hbar c)$ is the fine-structure
constant~\cite{MohrTaylor2018}.
Its current CODATA value is
$\Rinf = \SI{1.0973731568160e7}{\per\meter}$ with a relative uncertainty
of $\SI{1.9e-12}{}$.

\subsection{Reformulation in vortex-core variables}

\begin{lemma}[Rydberg identity in core variables]
\label{lem:rydberg}
Given Postulate~\ref{post:compton}, the Rydberg constant satisfies
\begin{equation}
  \Rinf = \frac{\vc^2\,\omC}{\pi c^3}
         = \frac{\vc^3}{\pi\,\rc\,c^3}.
  \label{eq:rydberg_core}
\end{equation}
\end{lemma}

\begin{proof}
Substitute $\rc = \vc/\omC$ into the right-hand side of
\eqref{eq:rydberg_core}:
\begin{equation}
  \frac{\vc^3}{\pi\,(\vc/\omC)\,c^3}
  = \frac{\vc^2\,\omC}{\pi c^3}
  = \frac{\vc^2\,m_e c^2}{\pi\,\hbar\,c^3}
  = \frac{\vc^2\,m_e}{\pi\,\hbar\,c}.
  \label{eq:rydberg_proof_step}
\end{equation}
Now set $\vc = \alpha c/2$ (which will be \emph{derived} rather
than assumed; see Sec.~\ref{sec:alpha}) and use
$\omC = m_e c^2/\hbar$:
\begin{equation}
  \frac{(\alpha c/2)^2\,m_e}{\pi\,\hbar\,c}
  = \frac{\alpha^2 m_e c}{4\pi\hbar}
  = \frac{\alpha^2 m_e c}{2h}
  = \Rinf,
  \label{eq:rydberg_id_check}
\end{equation}
which is the standard definition~\eqref{eq:Rydberg_standard}.
The identity holds exactly.
\end{proof}

\begin{remark}
Lemma~\ref{lem:rydberg} can be read in either direction.
If $\vc$ and $\rc$ are already known via $\alpha$, the identity is
algebraically guaranteed (a reformulation).
If $\vc$ is \emph{unknown} but linked to $\rc$ by
Postulate~\ref{post:compton}, then \eqref{eq:rydberg_core} becomes
a constraint that uniquely determines $\vc$ from $\Rinf$.
It is the second reading that is used in what follows.
\end{remark}

% =============================================================
\section{Derivation of the Characteristic Speed}
\label{sec:derivation}
% =============================================================

Substituting Postulate~\ref{post:compton} into Lemma~\ref{lem:rydberg}
eliminates $\rc$:
\begin{equation}
  \Rinf = \frac{\vc^2\,\omC}{\pi c^3}
  \qquad\Longrightarrow\qquad
  \vc^2 = \frac{\pi c^3 \Rinf}{\omC}
         = \frac{\pi\,\hbar\,c\,\Rinf}{m_e}.
  \label{eq:vc_solve}
\end{equation}
Taking the positive root,

\begin{proposition}[Characteristic circulation speed]
\label{prop:vc}
Given Postulate~\ref{post:compton} and the measured value of $\Rinf$,
\begin{equation}
  \boxed{
    \vc = \sqrt{\frac{\pi\,\hbar\,c\,\Rinf}{m_e}}
  }
  \label{eq:vc}
\end{equation}
is the unique positive solution for the characteristic circulation
speed.
The inputs are the standard constants $\{m_e,\,c,\,\hbar\}$ and the
spectroscopic observable $\Rinf$; the fine-structure constant $\alpha$
is not used.
\end{proposition}

\subsection{Dimensional check}
\begin{align}
  \Bigl[\sqrt{\frac{\hbar\,c\,\Rinf}{m_e}}\Bigr]
  &= \sqrt{\frac{\mathrm{J{\cdot}s}\cdot(\mathrm{m/s})\cdot\mathrm{m}^{-1}}
                {\mathrm{kg}}}
  = \sqrt{\frac{\mathrm{kg{\cdot}m^2{\cdot}s^{-2}{\cdot}s\cdot m{\cdot}m^{-1}}}
               {\mathrm{kg{\cdot}s}}}
  = \sqrt{\mathrm{m^2\,s^{-2}}}
  = \mathrm{m\,s^{-1}}. \notag
\end{align}

\subsection{Numerical evaluation}
Using CODATA 2018 values~\cite{MohrTaylor2018}:
\begin{align}
  m_e   &= \SI{9.1093837015e-31}{\kilogram}, \notag\\
  c     &= \SI{299792458}{\meter\per\second}, \notag\\
  \hbar &= \SI{1.054571817e-34}{\joule\second}, \notag\\
  \Rinf &= \SI{1.0973731568160e7}{\per\meter}, \notag
\end{align}
equation~\eqref{eq:vc} yields
\begin{equation}
  \vc = \SI{1.093\,845\,631\,5e6}{\meter\per\second}.
  \label{eq:vc_numerical}
\end{equation}

% =============================================================
\section{Derived Scales}
\label{sec:derived}
% =============================================================

With $\vc$ fixed by Proposition~\ref{prop:vc} and
Postulate~\ref{post:compton}, the following mechanical scales
are algebraic consequences.
No further assumptions are introduced.

% ------- 4.1 Core radius ------------------------------------
\subsection{Core radius}
\label{sec:rc}

From \eqref{eq:vc_rc_link},
\begin{equation}
  \rc = \frac{\vc}{\omC} = \frac{\vc\,\hbar}{m_e c^2}.
  \label{eq:rc}
\end{equation}
Numerically,
\begin{equation}
  \rc = \SI{1.408\,970\,16e-15}{\meter}.
  \label{eq:rc_numerical}
\end{equation}
This is half the classical electron radius, $\rc = r_e/2$, where
$r_e = \alpha\hbar/(m_e c)$.
The connection to $r_e$ follows from the derived value of $\alpha$
in Sec.~\ref{sec:alpha}.

% ------- 4.2 Fine-structure constant as output ---------------
\subsection{Fine-structure constant as a derived quantity}
\label{sec:alpha}

\begin{proposition}[Fine-structure constant]
\label{prop:alpha}
The dimensionless ratio
\begin{equation}
  \tilde\alpha \equiv \frac{2\vc}{c}
  \label{eq:alpha_derived}
\end{equation}
reproduces the CODATA fine-structure constant to within the precision
of the input data.
\end{proposition}

\begin{proof}
Substitute \eqref{eq:vc} into \eqref{eq:alpha_derived}:
\begin{equation}
  \tilde\alpha = \frac{2}{c}\sqrt{\frac{\pi\hbar c\Rinf}{m_e}}
  = \frac{2\sqrt{\pi\hbar\Rinf}}{\sqrt{m_e c}}.
\end{equation}
Using the standard identity $\Rinf = \alpha^2 m_e c/(2h)$
and $h = 2\pi\hbar$:
\begin{equation}
  \tilde\alpha = \frac{2\sqrt{\pi\hbar\cdot\alpha^2 m_e c/(2h)}}{\sqrt{m_e c}}
  = \frac{2\sqrt{\pi\hbar\alpha^2/(4\pi\hbar)}}{\,1\,}
  = 2\cdot\frac{\alpha}{2} = \alpha.
\end{equation}
The proof is algebraic; the numerical match to CODATA is a consistency
check, not an independent measurement.
\end{proof}

Numerically, $\tilde\alpha = \num{7.297\,352\,567\,1e-3}$, compared
with the CODATA value
$\alpha = \num{7.297\,352\,569\,3e-3}$;
the fractional difference is $\SI{3e-10}{}$, well within the
precision of $\Rinf$ itself.

\begin{remark}
Proposition~\ref{prop:alpha} reverses the usual logical order.
In standard treatments, $\alpha$ is input and $\Rinf$ is derived.
Here, $\Rinf$ is input and $\alpha$ is derived.
The two formulations are mathematically equivalent but carry different
epistemic weight: Proposition~\ref{prop:alpha} demonstrates that
a single spectroscopic measurement is sufficient to pin down $\alpha$,
given Postulate~\ref{post:compton}.
\end{remark}

% ------- 4.3 Maximum vortex tension --------------------------
\subsection{Maximum vortex tension}
\label{sec:Fmax}

The energy stored in the core over the scale $\rc$ defines a
characteristic force,
\begin{equation}
  \Fmax \equiv \frac{\hbar\omC}{2\rc}
              = \frac{m_e c^2}{2\rc}.
  \label{eq:Fmax_def}
\end{equation}
Using \eqref{eq:rc} and $\alpha = 2\vc/c$:

\begin{corollary}[Maximum tension in fundamental constants]
\label{cor:Fmax}
\begin{equation}
  \Fmax = \frac{m_e^2 c^3}{\alpha\hbar}
        = \frac{\hbar\omC^2}{2\vc},
  \label{eq:Fmax}
\end{equation}
dimensionally $[\Fmax] = \SI{}{\newton}$.
\end{corollary}

The derivation is three lines:
\begin{align}
  \Fmax = \frac{\hbar\omC}{2\rc}
        = \frac{\hbar\omC}{2\vc/\omC}
        = \frac{\hbar\omC^2}{2\vc}
        = \frac{\hbar(m_e c^2/\hbar)^2}{2\vc}
        = \frac{m_e^2 c^4}{2\hbar\vc}
        = \frac{m_e^2 c^3}{\alpha\hbar}.
\end{align}
Numerically, $\Fmax = \SI{29.053\,51}{\newton}$.

An equivalent algebraic form (derivable by the same substitutions) is
\begin{equation}
  \Fmax = \frac{\alpha\hbar c}{4\rc^2}
         = \frac{e^2}{16\pi\varepsilon_0\rc^2},
  \label{eq:Fmax_em}
\end{equation}
showing that $\Fmax$ is the Coulomb force between two elementary
charges separated by $2\rc$ times a numerical factor.
This is an internal consistency check, not an independent derivation.

\subsubsection*{Relation to the gravitational force scale}

In general relativity the combination $c^4/(4G)$ defines the maximum
force a classical process can transmit without forming a
horizon~\cite{Gibbons2002,Schiller2006}:
\begin{equation}
  \Fmax^{\mathrm{gr}} = \frac{c^4}{4G}.
  \label{eq:Fgr}
\end{equation}
The ratio of the two tension scales reduces to
\begin{equation}
  \frac{\Fmax}{\Fmax^{\mathrm{gr}}}
  = \frac{4\alpha_g}{\alpha},
  \label{eq:tension_ratio}
\end{equation}
where $\alpha_g = Gm_e^2/(\hbar c) \approx \num{1.752e-45}$ is the
gravitational fine-structure constant of the electron.
This is a dimensionless scale-bridge identity; its derivation uses
only the definitions of $\Fmax$, $\Fmax^{\mathrm{gr}}$, and $\alpha_g$,
and introduces no new parameters.

% ------- 4.4 Core mass density --------------------------------
\subsection{Core mass density}
\label{sec:rhocore}

The cylindrical Laplace pressure across a vortex tube of radius $\rc$
and line tension $\mathcal{T} = \Fmax/(2\pi\rc)$ is
$P = \mathcal{T}/\rc$~\cite{Saffman1992,Donnelly1991}.
Setting this equal to the kinetic pressure $\frac{1}{2}\rhocore\vc^2$
of the core fluid:
\begin{equation}
  \frac{1}{2}\rhocore\vc^2 = \frac{\Fmax}{2\pi\rc^2}.
  \label{eq:pressure_balance}
\end{equation}
Solving for $\rhocore$ and substituting \eqref{eq:Fmax_def}:
\begin{equation}
  \rhocore = \frac{\Fmax}{\pi\rc^2\vc^2}
           = \frac{\hbar\omC}{2\pi\rc^3\vc^2}
           = \frac{m_e c^2}{2\pi\vc^2\rc^3}.
  \label{eq:rhocore}
\end{equation}

\begin{corollary}[Core density]
\label{cor:rho}
\begin{equation}
  \boxed{
    \rhocore = \frac{m_e c^2}{2\pi\,\vc^2\,\rc^3}
  }
  \label{eq:rhocore_boxed}
\end{equation}
with $[\rhocore] = \SI{}{\kilogram\per\cubic\meter}$.
\end{corollary}

Dimensional check:
$[m_e c^2 / (\vc^2 \rc^3)]
= \mathrm{(kg\,m^2\,s^{-2})\,(m^2\,s^{-2})^{-1}\,m^{-3}}
= \mathrm{kg\,m^{-3}}$ \checkmark.

Numerically, $\rhocore = \SI{3.893\,4e18}{\kilogram\per\cubic\meter}$.

\begin{remark}
The derivation of $\rhocore$ uses only $\{m_e, c, \vc, \rc\}$,
where both $\vc$ and $\rc$ were obtained without circular reference
to $\rhocore$ itself.
The Laplace-pressure balance~\eqref{eq:pressure_balance} is the
physical content; the result~\eqref{eq:rhocore_boxed} is its
consequence.
The ratio $\rhocore/\rho_{\Lambda}$, where
$\rho_{\Lambda} = \Lambda c^2/(8\pi G) \approx \SI{5.84e-27}{\kilogram\per\cubic\meter}$
is the cosmological vacuum energy density, is approximately
$\num{6.7e44}$;
this large ratio has no mechanistic explanation within the present
framework.
\end{remark}

% ------- 4.5 Circulation quantum ----------------------------
\subsection{Circulation quantum}
\label{sec:Gamma0}

The circulation around the vortex core of radius $\rc$ with tangential
speed $\vc$ defines the vortex's circulation:
\begin{equation}
  \Gz \equiv 2\pi\,\rc\,\vc = \frac{2\pi\vc^2}{\omC}.
  \label{eq:Gamma0}
\end{equation}
Numerically, $\Gz = \SI{9.683\,619e-9}{\meter\squared\per\second}$.

For comparison, the Onsager-Feynman circulation quantum in a
superfluid of particle mass $m_e$ is
$\kappa = h/m_e \approx \SI{7.274e-4}{\meter\squared\per\second}$~\cite{Onsager1949,Feynman1955}.
The ratio $\Gz/\kappa = \alpha^2/4 \approx \num{1.33e-5}$, showing
that the present circulation quantum is smaller than the
Onsager-Feynman quantum by the square of the fine-structure constant
divided by four.

% ------- 4.6 Geometric gate identity -------------------------
\subsection{Geometric gate identity}
\label{sec:gate}

The ratio of the electron Compton wavelength $\lc = h/(m_e c)$ to
$\pi\rc$ is

\begin{lemma}[Geometric gate]
\label{lem:gate}
\begin{equation}
  \frac{\lc}{\pi\rc} = \frac{4}{\alpha}.
  \label{eq:gate}
\end{equation}
\end{lemma}

\begin{proof}
$\lc = h/(m_e c) = 2\pi\hbar/(m_e c)$.
Substituting $\rc = \alpha\hbar/(2m_e c)$ (which follows from $\rc = \vc/\omC$
and $\alpha = 2\vc/c$):
\begin{equation}
  \frac{\lc}{\pi\rc}
  = \frac{2\pi\hbar/(m_e c)}{\pi\,\alpha\hbar/(2m_e c)}
  = \frac{4}{\alpha}.
\end{equation}
\end{proof}

Numerically, $4/\alpha = \lc/(\pi\rc) = \num{548.144}$.
This dimensionless ratio serves as a natural amplification factor in
mass-energy calculations for the vortex core~\cite{CantarellaKusner2002}.

% =============================================================
\section{Derivation Chain Summary}
\label{sec:chain}
% =============================================================

The complete non-circular derivation chain is:
\begin{equation}
  \underbrace{\omC,\;\Rinf,\;m_e,\;c,\;\hbar}_{\text{inputs}}
  \xrightarrow{\text{Post.\;\ref{post:compton}}}
  \vc
  \rightarrow
  \rc
  \rightarrow
  \tilde\alpha
  \rightarrow
  \Fmax
  \rightarrow
  \rhocore
  \rightarrow
  \Gz.
  \label{eq:chain}
\end{equation}
Each arrow represents an algebraic step.
No quantity downstream uses any quantity upstream as a free parameter.

% =============================================================
\section{Status Table}
\label{sec:status}
% =============================================================

\begin{table}[h]
\centering
\caption{Epistemic status of all quantities in the derivation.
``Derived'' means the quantity follows algebraically from the stated
inputs and Postulate~\ref{post:compton}.
``Calibrated'' means it is supplied as a measured input.
``Open'' means it is not determined within the present framework.}
\label{tab:status}
\begin{tabular}{llll}
\toprule
Quantity & Symbol & Status & Source \\
\midrule
Compton frequency    & $\omC = m_e c^2/\hbar$   & Derived  & standard defs.~\cite{MohrTaylor2018} \\
Core rotation rate   & $\Omega_{\rm core}=\omC$ & Postulate & Post.~\ref{post:compton} \\
Rydberg constant     & $\Rinf$                  & Calibrated (input) & spectroscopy~\cite{MohrTaylor2018} \\
Circulation speed    & $\vc$                    & Derived  & Prop.~\ref{prop:vc} \\
Core radius          & $\rc$                    & Derived  & \eqref{eq:rc} \\
Fine-structure const.& $\tilde\alpha = 2\vc/c$  & Derived  & Prop.~\ref{prop:alpha} \\
Maximum tension      & $\Fmax$                  & Derived  & Cor.~\ref{cor:Fmax} \\
Core density         & $\rhocore$               & Derived  & Cor.~\ref{cor:rho} \\
Circulation quantum  & $\Gz$                    & Derived  & \eqref{eq:Gamma0} \\
Geometric gate       & $\lc/(\pi\rc) = 4/\alpha$& Derived  & Lemma~\ref{lem:gate} \\
Background density   & $\rho_f$                 & Open     & not determined \\
Topological state    & knot type $K$            & Open     & not determined \\
\bottomrule
\end{tabular}
\end{table}

The one quantity that remains open is the background fluid density
$\rho_f$ (the density of the vortex medium far from the core).
The present framework does not determine it; it requires an additional
physical constraint, such as a measurement of the vorticity-induced
refractive index $\Delta n = \rho_f|\boldsymbol{\omega}|^2/c^2$ in a
rotating medium at known angular velocity
(see Sec.~\ref{sec:falsifiers}).

% =============================================================
\section{Falsifiable Predictions}
\label{sec:falsifiers}
% =============================================================

The following three predictions hold for any model sharing
Postulate~\ref{post:compton} and Proposition~\ref{prop:vc}.
They involve no free parameters.

\medskip
\noindent\textbf{(i) Core-radius scaling.}
The core radius satisfies $\rc = \vc/\omC$ precisely.
Any independent experimental measurement of $\rc$ inconsistent with
$\rc = \SI{1.408\,970\,16e-15}{\meter}$ falsifies
Postulate~\ref{post:compton}.
High-precision electron scattering or Lamb-shift
measurements~\cite{Karshenboim2005} that detect an anomalous
short-range component of the Coulomb potential at the femtometre scale
provide the most direct test.

\medskip
\noindent\textbf{(ii) Tension-to-gravity ratio.}
Equation~\eqref{eq:tension_ratio} predicts
\begin{equation}
  \frac{\Fmax}{\Fmax^{\mathrm{gr}}} = \frac{4\alpha_g}{\alpha}
  \approx \num{9.60e-43}
  \label{eq:pred2}
\end{equation}
without any fitted parameters.
This can be tested in analogue-gravity experiments in Bose-Einstein
condensates~\cite{Barcelo2005} by comparing the maximum phonon emission
rate (analogue of $\Fmax$) to the Planck-scale force analogue of the
system.

\medskip
\noindent\textbf{(iii) Refractive-index constraint on background density.}
The model predicts that a uniformly rotating medium of background density
$\rho_f$ at angular velocity $|\boldsymbol{\omega}|$ produces a
refractive-index shift
\begin{equation}
  \Delta n = \frac{\rho_f\,|\boldsymbol{\omega}|^2}{c^2}.
  \label{eq:refraction}
\end{equation}
This can be inverted to determine $\rho_f$ from a single interferometric
measurement:
\begin{equation}
  \rho_f = \frac{\Delta n\,c^2}{|\boldsymbol{\omega}|^2}.
  \label{eq:rho_f_measurement}
\end{equation}
At a rotation rate $|\boldsymbol{\omega}| \approx \SI{1.1e4}{\per\second}$
(accessible in cold-atom or SQUID experiments) and interferometric
precision $\Delta n \sim 10^{-15}$, equation~\eqref{eq:rho_f_measurement}
resolves $\rho_f$ to an uncertainty of order $\SI{1e-7}{\kilogram\per\cubic\meter}$.
This is the only known non-circular route to the background density
within the vortex-core framework.

% =============================================================
\section{Numerical Summary}
\label{sec:numerical}
% =============================================================

\begin{table}[h]
\centering
\caption{Numerical values of all derived quantities, evaluated with
CODATA 2018 inputs~\cite{MohrTaylor2018}.
The last column gives the relative deviation from the corresponding
CODATA value where applicable.}
\label{tab:numbers}
\begin{tabular}{llll}
\toprule
Quantity & Symbol & Value & $\delta/\delta_{\rm CODATA}$ \\
\midrule
Circulation speed   & $\vc$         & \SI{1.093\,845\,63e6}{\meter\per\second}          & $\sim 10^{-9}$ \\
Core radius         & $\rc$         & \SI{1.408\,970\,16e-15}{\meter}                   & $\sim 10^{-9}$ \\
Fine-structure const.& $\tilde\alpha$ & \num{7.297\,352\,567e-3}                         & $3\times 10^{-10}$ \\
Max.\ tension       & $\Fmax$       & \SI{29.053\,51}{\newton}                          & — \\
Core density        & $\rhocore$    & \SI{3.893\,4e18}{\kilogram\per\cubic\meter}        & — \\
Circ.\ quantum      & $\Gz$         & \SI{9.683\,619e-9}{\meter\squared\per\second}     & — \\
Geometric gate      & $4/\alpha$    & \num{548.14}                                      & $3\times 10^{-10}$ \\
\bottomrule
\end{tabular}
\end{table}

% =============================================================
\section{Discussion}
\label{sec:discussion}
% =============================================================

The main result of this paper is that a single structural postulate
(Postulate~\ref{post:compton}) combined with one spectroscopic
observable ($\Rinf$) is sufficient to fix all mechanical scales of a
localized vortex-core model for the electron.
In particular, the fine-structure constant emerges as a derived ratio
rather than an independent input, and the derivation is free of the
circular arguments that appear when $\alpha$ is used to define both
the circulation speed and the core radius.

The derivation is purely algebraic: no differential equation is solved,
no numerical optimization is performed, and no particle mass other than
$m_e$ is used.
The result places all commonly used vortex-core parameters on a single
logical footing, with explicit status labels (Table~\ref{tab:status}).

The one genuinely open parameter, the background fluid density $\rho_f$,
is not determined by the present framework.
Equation~\eqref{eq:rho_f_measurement} shows how it could in principle
be measured.

This algebraic structure arose in the context of a toy model for
topological vortex excitations of the vacuum medium, but the
derivations in Secs.~\ref{sec:derivation}--\ref{sec:derived} are
entirely model-independent: they hold for any vortex-core description
that satisfies Postulate~\ref{post:compton}.

% =============================================================
\section{Conclusion}
\label{sec:conclusion}
% =============================================================

We have shown that the characteristic circulation speed
$\vc = (\pi\hbar c\Rinf/m_e)^{1/2}$,
the core radius $\rc = \vc/\omC$,
the fine-structure constant $\tilde\alpha = 2\vc/c$,
the maximum vortex tension $\Fmax = \hbar\omC/(2\rc)$,
the core density $\rhocore = m_ec^2/(2\pi\vc^2\rc^3)$,
and the circulation quantum $\Gz = 2\pi\vc^2/\omC$
follow from a single postulate about the core rotation rate
and the spectroscopic Rydberg constant.
The derivation is non-circular, dimensionally explicit, and agrees
with CODATA values to ten or more significant figures.
Three falsifiable predictions are stated.

% =============================================================
\begin{acknowledgments}
The author thanks multiple independent review processes for
identifying and correcting several earlier dimensional and
circularity errors.
\end{acknowledgments}

% =============================================================
\begin{thebibliography}{99}

\bibitem{Helmholtz1867}
H.~von Helmholtz,
On the integrals of the hydrodynamical equations which express
vortex-motion,
\textit{Philosophical Magazine} \textbf{33}, 485 (1867).

\bibitem{Kelvin1867}
Lord Kelvin (W.~Thomson),
On vortex atoms,
\textit{Proceedings of the Royal Society of Edinburgh} \textbf{6},
94 (1867).

\bibitem{MohrTaylor2018}
P.~J. Mohr, D.~B. Newell, B.~N. Taylor, and E.~Tiesinga,
CODATA recommended values of the fundamental physical constants: 2018,
\textit{Reviews of Modern Physics} \textbf{93}, 025010 (2021).
\doi{10.1103/RevModPhys.93.025010}

\bibitem{Donnelly1991}
R.~J. Donnelly,
\textit{Quantized Vortices in Helium II},
Cambridge University Press, 1991.

\bibitem{Pethick2008}
C.~J. Pethick and H.~Smith,
\textit{Bose–Einstein Condensation in Dilute Gases}, 2nd ed.,
Cambridge University Press, 2008.
\doi{10.1017/CBO9780511802850}

\bibitem{FaddeevNiemi1997}
L.~D. Faddeev and A.~J. Niemi,
Stable knot-like structures in classical field theory,
\textit{Nature} \textbf{387}, 58 (1997).
\doi{10.1038/387058a0}

\bibitem{Barcelo2005}
C.~Barceló, S.~Liberati, and M.~Visser,
Analogue gravity,
\textit{Living Reviews in Relativity} \textbf{8}, 12 (2005).
\doi{10.12942/lrr-2005-12}

\bibitem{Volovik2003}
G.~E. Volovik,
\textit{The Universe in a Helium Droplet},
Oxford University Press, 2003.

\bibitem{Saffman1992}
P.~G. Saffman,
\textit{Vortex Dynamics},
Cambridge University Press, 1992.
\doi{10.1017/CBO9780511624063}

\bibitem{Onsager1949}
L.~Onsager,
Statistical hydrodynamics,
\textit{Nuovo Cimento Supplement} \textbf{6}, 279 (1949).
\doi{10.1007/BF02780991}

\bibitem{Feynman1955}
R.~P. Feynman,
Application of quantum mechanics to liquid helium,
in \textit{Progress in Low Temperature Physics}, vol.~1,
ed.\ C.~J. Gorter, p.~17,
North-Holland, 1955.

\bibitem{Gibbons2002}
G.~W. Gibbons,
The maximum tension principle in general relativity,
\textit{Foundations of Physics} \textbf{32}, 1891 (2002).
\doi{10.1023/A:1022370717626}

\bibitem{Schiller2006}
C.~Schiller,
Simple derivation of minimum length, minimum dipole moment, and
lack of space-time continuity,
\textit{International Journal of Theoretical Physics} \textbf{45},
221 (2006).
\doi{10.1007/s10773-005-4835-2}

\bibitem{CantarellaKusner2002}
J.~Cantarella, R.~B. Kusner, and J.~M. Sullivan,
On the minimum ropelength of knots and links,
\textit{Inventiones Mathematicae} \textbf{150}, 257 (2002).
\doi{10.1007/s00222-002-0234-y}

\bibitem{Karshenboim2005}
S.~G. Karshenboim,
Precision physics of simple atoms: QED tests, nuclear structure, and
fundamental constants,
\textit{Physics Reports} \textbf{422}, 1 (2005).
\doi{10.1016/j.physrep.2005.08.008}

\bibitem{Bethe1947}
H.~A. Bethe,
The electromagnetic shift of energy levels,
\textit{Physical Review} \textbf{72}, 339 (1947).
\doi{10.1103/PhysRev.72.339}

\bibitem{Batchelor1967}
G.~K. Batchelor,
\textit{An Introduction to Fluid Dynamics},
Cambridge University Press, 1967.
\doi{10.1017/CBO9780511800955}

\end{thebibliography}

\end{document}
